\section{Introduction}
\label{sec:introduction}

For an integer $k\ge2$, an integer $a$ is \emph{$k$-power-free} if no prime
power $p^k$ divides $a$. This convention includes negative integers and
excludes zero. Given $f\in\Z[x]$, write
\[
 \rho_f(q)=\#\{a\in\Z/q\Z:f(a)\equiv0\pmod q\},\qquad
 S_{f,k}(X)=\#\{1\le n\le X:f(n)\text{ is $k$-power-free}\}.
\]
The necessary local condition is
\begin{equation}
 \rho_f(p^k)<p^k\quad\text{for every prime }p.
 \label{eq:local-admissibility}
\end{equation}
It says that no prime $k$th power divides every value of $f$. The natural
candidate for the density is the product of the proportions that survive
these separate congruence conditions.

\begin{theorem}\label{thm:density}
Let $f\in\Z[x]$ be irreducible over $\Q$, of degree $d$ with
$4\le d\le8$, and put $k=d-2$. If \eqref{eq:local-admissibility} holds,
then
\[
 S_{f,k}(X)=c_{f,k}X+o_f(X),\qquad
 c_{f,k}=\prod_p\left(1-\frac{\rho_f(p^k)}{p^k}\right)>0.
\]
No primitivity or sign condition on the leading coefficient is required.
\end{theorem}

For quartics this is a squarefree-value density theorem. In degrees
$5,6,7,8$, the respective excluded powers are $3,4,5,6$.
The statement concerns positive integer inputs and imposes no bound on
the height of the fixed polynomial's coefficients. The error term is not
asserted to be uniform as $f$ varies.

In particular, the quartics $x^4+2$ and $x^4+1$ have positive squarefree-value
density. The first is Eisenstein at $2$; the second is the eighth cyclotomic
polynomial. For $x^4+2$, its value at zero excludes a fixed odd prime square,
and its value at one excludes a fixed square of $2$. The value at zero
of $x^4+1$ is one, so its local condition is immediate.

\subsection{Earlier results and the higher-degree consequence}

The power-free-value problem has a long history. Ricci obtained
the asymptotic in the range $k\ge d$ \cite{Ricci1933}.
For $d\ge3$, Erd\H{o}s proved infinitude for exponent $d-1$, and Hooley obtained
the corresponding asymptotic \cite{Erdos1953,Hooley1967}; see also
the historical discussion in \cite[Introduction]{HeathBrown2013Powerfree}.
The cubic squarefree case therefore belongs to this classical range.
Already in 1953, Erd\H{o}s singled out the unresolved squarefreeness
of $n^4+2$ \cite[p.~425]{Erdos1953}. In his 1965 survey he returned
to this example when describing the obstacle at exponent $d-2$
\cite[Section~6, p.~219]{Erdos1965}.

Nair obtained the asymptotic for
$k\ge(\sqrt2-\tfrac12)d$ \cite{Nair1976}, with further developments
by Nair and Huxley--Nair \cite{Nair1979,HuxleyNair1980}.
The 1976 bound reaches $k=d-2$ for $d\ge24$.
Heath-Brown's affine determinant method gave the range
$k\ge(3d+2)/4$ \cite[Theorem~16]{HeathBrown2006Counting}, reaching
$k=d-2$ for $d\ge10$. Combining that method with Salberger's
global determinant estimates, Browning improved the range to
$k\ge(3d+1)/4$ \cite{Browning2011PowerFree}, which reaches
$k=d-2$ for $d\ge9$.
We use the formulation and published reproof in
Xiao \cite[Theorem~9.1 and Section~9]{Xiao2017PowerFree}:
the input interval is $1\le n\le X$ and the constant is the
Euler product over all residue classes modulo $p^k$.

The shape of the polynomial can permit stronger estimates.
Heath-Brown proved an asymptotic with a power-saving error for
irreducible binomials $x^d+c$ when $k\ge(5d+3)/9$
\cite[Theorem~1]{HeathBrown2013Powerfree}. At $k=d-2$ this covers
such binomials in degrees $d\ge6$, while imposing the special
binomial form. Reuss obtained a power-saving error for the general
degree-$d$ polynomial at exponent $d-1$
\cite[Theorem~2]{Reuss2015PowerFree}. The quartic theorem above
concerns squarefree values at positive integer arguments for each
fixed admissible polynomial.

Other important directions involve different hypotheses or
quantifiers. Granville showed that the $abc$ conjecture yields
the predicted squarefree-value density for separable integer
polynomials with no fixed prime-square divisor
\cite[Theorem~1]{granville1998abc}.
Greaves and Helfgott studied squarefree values of binary forms;
Helfgott's Theorems~5.1--5.2 distinguish univariate cubics from
homogeneous forms in two variables of degree at most six
\cite{Greaves1992,Helfgott2004}. A density in two integer variables
does not by itself give the density on the fixed line with second
coordinate one. Recent results of Browning--Shparlinski and Sofos
study the discrepancy from the predicted squarefree count on
average over polynomial coefficient families
\cite[Theorems~1.1--1.2]{BrowningShparlinski2024},
\cite[Theorem~1.1]{Sofos2026}; those averaged statements
do not imply a theorem for every fixed polynomial.

Carella's preprint claims squarefree-value asymptotics for both
$n^4+1$ and $n^4+2$ \cite[Theorems~1.1--1.2]{carella2023squarefree}.
Zapata Ceballos--Jalalvand claim positive density for monic
irreducible degree-$2q$ polynomials with squarefree fixed divisor
whose root field contains a degree-$q$ Galois subfield, where $q$
is prime; this includes $n^4+1$
\cite[Theorem~1.1]{zapataceballosjalalvand2026squarefree}.
We do not use either claim as a
theorem input. Our large-prime estimate and sieve transfer are
proved independently below.

The new argument in this paper concerns degrees four through eight.
For higher degrees, Browning's theorem supplies a direct consequence.

\begin{corollary}\label{cor:all-degrees}
Let $f\in\Z[x]$ be irreducible over $\Q$ of degree $d\ge4$, and
put $k=d-2$. If $f$ has no fixed prime $k$th-power divisor, then
\[
 S_{f,k}(X)=c_{f,k}X+o_f(X),\qquad
 c_{f,k}=\prod_p\left(1-\frac{\rho_f(p^k)}{p^k}\right)>0.
\]
\end{corollary}

\begin{proof}
For $4\le d\le8$ use Theorem~\ref{thm:density}. For $d\ge9$,
\[
 d-2\ge\frac{3d+1}{4},
\]
so Browning's theorem, in the cited form of Xiao, applies to the
primitive positive-leading part of $f$. Lemma~\ref{lem:density-content}
transfers its density to $f$, retaining the original local factors.
\end{proof}

The factorwise large-prime estimates also give a fixed-product
consequence. Here a nonzero polynomial is \emph{separable over $\Q$}
if each of its nonconstant irreducible factors occurs with multiplicity one.

\begin{corollary}[Separable products and simultaneous values]
\label{cor:separable-products}
Fix an integer $k\ge2$. Let $F\in\Z[x]\setminus\{0\}$ be separable over
$\Q$, with every nonconstant irreducible factor of degree at most $k+2$.
If $\rho_F(p^k)<p^k$ for every prime $p$, then
\[
 S_{F,k}(N)=c_{F,k}N+o_{F,k}(N),\qquad
 c_{F,k}=\prod_p\left(1-\frac{\rho_F(p^k)}{p^k}\right)>0.
\]
More generally, let $\mathbf F=(F_1,\ldots,F_r)$, with $r\ge1$, be a fixed
family of nonzero integer polynomials, each separable over $\Q$ with all
nonconstant irreducible factors of degree at most $k+2$. Factors may be
shared by different $F_i$. Put
\[
 \Omega_p=\{a\in\Z/p^k\Z:p^k\mid F_i(a)\text{ for at least one }i\}.
\]
Let $S_{\mathbf F,k}(N)$ count the integers $1\le n\le N$ for which every
$F_i(n)$ is $k$-power-free. If $\#\Omega_p<p^k$ for every prime $p$, then
\[
 S_{\mathbf F,k}(N)=c_{\mathbf F,k}N+o_{\mathbf F,k}(N),\qquad
 c_{\mathbf F,k}=\prod_p\left(1-\frac{\#\Omega_p}{p^k}\right)>0.
\]
\end{corollary}

The first assertion requires the value of the entire product $F$ to be
$k$-power-free. The second imposes power-freeness separately on the
$F_i(n)$, not on their product, and uses the joint local condition
$\Omega_p$ rather than a product of individual local factors.
All polynomials, their number, and $k$ are fixed; no uniform or
power-saving error is asserted here. The proof is at the end of
Section~\ref{sec:sieve}.

Our method combines number-field factorization with approximating
determinants, following the framework developed by Reuss for
power-free polynomial values \cite[Sections~3--6]{Reuss2015PowerFree}.
Heath-Brown's determinant arguments for
squarefree values of $n^2+1$ and for power-free binomial values
provide related precedents
\cite{HeathBrown2012Squarefree,HeathBrown2013Powerfree}.
The real determinant method originates in work of Bombieri--Pila;
Heath-Brown developed the prime-adic and approximating forms used
in this line of arguments, and Salberger developed their local
Hilbert and global congruence aspects
\cite{BombieriPila1989,HeathBrown2002Density,HeathBrown2009Sums,
Salberger2007,Salberger2023}.
The uniform affine counting theorem, the explicit weighted Hilbert
calculations and the quartic surface and curve arguments needed here
are proved below.

\subsection{The large-prime problem}

The finite-prime part of the problem follows from the Chinese remainder
theorem. For primes outside a fixed finite set, separability and Hensel
lifting give $\rho_f(p^k)\le d$. Consequently, for fixed large $Y$, primes
$Y<p\le X$ remove at most
\[
 O_f\bigl(XY^{1-k}+\pi(X)\bigr)
\]
integers in $(X,2X]$. The remaining task is to control primes larger than
the length of the input interval. We prove the following estimate.

\begin{proposition}\label{prop:large-prime-tail}
Let $f\in\Z[x]$ be primitive and irreducible, with positive leading
coefficient and degree $4\le d\le8$. Put $k=d-2$. There exists
$\delta=\delta_f>0$ such that
\[
 \#\{n\in\Z:X<n\le2X,\ p^k\mid f(n)
          \text{ for some prime }p>X\}
 \ll_f X^{1-\delta}.
\]
\end{proposition}

The proposition does not require local admissibility. In
Section~\ref{sec:sieve}, we prove explicitly that it implies
Theorem~\ref{thm:density}, including the stated sign and content
generality. The density argument first holds the finite-prime cutoff
fixed while $X$ tends to infinity, and then lets the cutoff tend to
infinity. This order gives the exact Euler product without an assertion
of a power-saving error in the density asymptotic.

\subsection{Proof strategy}

Fix a dyadic prime interval $P\le p\le2P$, and choose one such prime
for each input being counted. Write $P=X^\eta$ and $b=d-k\eta$.
The elementary relation $f(n)=yp^k$ puts the chosen triples on the
surface $f(x)=yz^k$, with coordinate sizes $X,X^b,X^\eta$.
A weighted determinant argument on this surface suffices when $\eta$
is large. The harder range requires more information from the
factorization of $f$.

Let $\theta$ be a root of $f$ and let $K=\Q(\theta)$. Ideal
factorization and the unit theorem produce integral elements
$\alpha,\beta\in K$ satisfying
\[
 \alpha^k\beta=\mu(n-\theta)
\]
for a fixed nonzero integer $\mu$. Their conjugates are balanced:
the sizes of $\alpha$ and $\beta$ at each embedding are respectively
comparable to $P^{1/d}$ and $X^{b/d}$. Integer coordinates in a fixed
integral basis therefore give two boxes, while their projective
directions satisfy explicit polynomial equations. Near a balanced
point, and after choosing one of finitely many local branches, either
direction determines the other analytically once $1/n$ is specified.

We partition these directions into small paired boxes. An
archimedean determinant gives a polynomial equation on the points
in each pair. Components defined by an equation in only one of the
two coordinate groups receive a second cut. Explicit complete
intersection calculations measure the independent monomials available
on the resulting sets. The distinction between a mixed equation and a
pure equation is important: the latter also reduces the number of
independent Taylor monomials in the local determinant.

Two further steps turn these cuts into a count. First, a reduced basis
of the integer lattice associated with a direction box makes its
points smaller in new integral coordinates. An exceptionally short
lattice vector also limits the number of boxes in which it can occur.
Second, we apply the general affine counting theorem proved in
Section~\ref{sec:affine-counting}, using lower bounds
for weighted spaces of polynomial functions. Its proof repeatedly
uses residue classes at small auxiliary primes. Its auxiliary
varieties depend on earlier choices of primes; a random deletion
argument controls the primes at which their equations fail to give
smooth local coordinates. This permits a uniform count through
successive dimension drops.

For quartics, surfaces and curves need sharper geometric bounds.
On a surface, the difficult projections have degree one or two.
The quadric case supplies an additional field degree. On a plane,
a deficient bound would force either a polynomial of degree at most
four on a general line or a degree-two rational map on the line of
directions to have four distinct critical values. On a curve, the limiting
case is a polynomial of degree five with four prescribed critical
values. Up to a normalization of two critical points, only finitely
many such polynomials occur. Their integer values in the input interval
form one set of size $O_f(X^{1/5})$, which is removed once before summing
over boxes. The remaining curve estimates and an exact finite choice
of parameters give a common power saving.

\subsection{Organization}

Figure~\ref{fig:proof-flow} shows how the principal estimates feed
into the density theorem. Its arrows indicate uses of established
modules; the order of the sections places the uniform counting
statement before its arithmetic applications.

\begin{figure}[ht]
\centering
\begin{tikzpicture}[
  node distance=8mm and 7mm,
  flowbox/.style={draw,rounded corners=2pt,align=center,text width=0.40\textwidth,
    inner sep=5pt,font=\small},
  arrow/.style={-{Latex[length=2mm]},semithick}]
\node[flowbox] (sieve) {Exact density from finite congruences\\
  Section~\ref{sec:sieve}};
\node[flowbox,right=of sieve] (tail) {Large prime powers $p^{d-2}\mid f(n)$\\
  Section~\ref{sec:tail-proof}};
\node[flowbox,below=of tail] (boxes) {Balanced field factors and paired boxes\\
  Sections~\ref{sec:arithmetic-boxes}--\ref{sec:lattice-boxes}};
\node[flowbox,left=of boxes] (hilbert) {Function spaces on surfaces and curves\\
  Sections~\ref{sec:hilbert-bounds}--\ref{sec:curves}};
\node[flowbox,below=of hilbert] (count) {Uniform affine counting\\
  Section~\ref{sec:affine-counting}};
\node[flowbox,right=of count] (parameters) {Strict parameter margins and summation\\
  Sections~\ref{sec:parameters}--\ref{sec:tail-proof}};
\draw[arrow] (boxes) -- (hilbert);
\draw[arrow] (hilbert) -- (count);
\draw[arrow] (count) -- (parameters);
\draw[arrow] (parameters.east) -- ++(5mm,0) |- (tail.east);
\draw[arrow] (tail) -- (sieve);
\end{tikzpicture}

\caption{The number-field route to the large-prime bound and the
exact density. The direct-surface estimate handles the complementary
upper prime range in Section~\ref{sec:tail-proof}.}
\label{fig:proof-flow}
\end{figure}

Section~\ref{sec:sieve} gives the density reduction, and
Section~\ref{sec:affine-counting} proves the general counting theorem.
Sections~\ref{sec:arithmetic-boxes}--\ref{sec:lattice-boxes} construct
the number-field boxes, determinant cuts and lattice coordinates.
Section~\ref{sec:hilbert-bounds} establishes weighted Hilbert bounds;
Sections~\ref{sec:quartic-geometry} and~\ref{sec:curves} sharpen them
for quartic surfaces and curves. Section~\ref{sec:parameters} selects
the parameters, and Section~\ref{sec:tail-proof} proves the
large-prime estimate and completes the density theorem.
Section~\ref{sec:cyclotomic} gives the explicit cyclotomic density,
retains a quantitative central-prime estimate, and explains the
special factorization, level equation and Pell alternative.
